\caption{Figure B. Training-set neighborhood diagnostics on historical ImageNet-A seed0 checkpoints. From left to right: mean positive Task0 cross-entropy increase; mean Task0 training accuracy (\%); mean raw empirical Fisher trace; and the mean-damage empirical tolerance radius. Radius is the joint Frobenius distance across 24 effective Q/V weight matrices from each method's own Task0 anchor. The same 16 unit directions are shared between methods and all positive radii; radius zero is one reused anchor. Loss and accuracy use all 1,259 Task0 training images with deterministic transforms and each method's fixed 20-class head; Fisher uses the same fixed 256 training images, true-label per-example gradients, without trace normalization, flooring or regularization. The first three panels show means across directions. Absolute accuracy retains the methods' different anchor accuracies. Positive damage is $[L_0(w_0+ru)-L_0(w_0)]_+$; its directional mean is $G(r)$. The tolerance radius $\hat r_\epsilon$ is the largest sampled radius such that $G(r_i)\leq\epsilon$ at every sampled radius up to it. Both methods have identical tolerance values at the five measured thresholds, shown with nested markers. Triangles mark right-censoring at the maximum scanned radius 20: the measured tolerance reaches the scan limit. All axes are linear; lines guide the eye between measured points. Only training means are plotted, with no test curves or Q95 series. These finite-direction, single-seed diagnostics are not worst-case guarantees.}
